\begin{figure}[htbp]
\centering
\begin{tikzpicture}[x=1cm,y=1cm]
\node[dissertation block,text width=2.75cm] (object) at (-5.65,1.35)
  {Початковий об’єкт\\$X_{mm}$};
\node[dissertation block,text width=2.75cm] (payload) at (-5.65,-0.65)
  {Приховане навантаження\\$P_A$};
\node[dissertation block,text width=3.00cm] (prepare) at (-2.10,0.35)
  {Підготовка об’єкта\\$E_A(\Theta_A)$};
\node[dissertation block,text width=2.75cm] (prepared) at (1.20,0.35)
  {Підготовлений об’єкт\\$X_{mm}^{\prime}$};
\node[dissertation block,text width=3.15cm] (conditions) at (4.65,0.35)
  {Умови прихованості\\$V_f=1$, $D\leq\varepsilon$};
\node[dissertation block,text width=3.30cm] (route) at (4.65,-1.75)
  {Маршрут $G$\\$r_j\in R_{web}$};
\node[dissertation block,text width=4.40cm] (observations) at (0.55,-3.20)
  {Статичні та подієві спостереження\\$X_s,X_v,X_f,X_b$};

\coordinate (inputbus) at (-4.05,0.35);
\draw[thick] (object.east) -- (object.east -| inputbus) -- (inputbus);
\draw[thick] (payload.east) -- (payload.east -| inputbus) -- (inputbus);
\fill[black] (inputbus) circle (1.3pt);
\draw[dissertation arrow] (inputbus) -- (prepare.west);
\draw[dissertation arrow] (prepare.east) -- (prepared.west);
\draw[dissertation arrow] (prepared.east) -- (conditions.west);
\draw[dissertation arrow] (conditions.south) -- (route.north);
\draw[dissertation arrow] (route.south) |- (observations.east);
\end{tikzpicture}
\caption{Причинна схема формування спостережень за умов прихованої доставки мультимедійного об’єкта}
\label{fig:attack_model}
\end{figure}
